\subsection{Denotational Semantics of \GRFRP}
\label{sec:semantics:grfrp:denot}

The denotational semantics of \GRFRP, following previous work~\cite{DBLP:journals/tcad/LeeS98},
associates to flows a timestamped stream representing their evolution over time: the timestamped
values in a signal stream represent its updates, while those of an event represent its occurrences.
This section reviews its construction.

\begin{table}[!ht]
    \centering
    \begin{tabularx}{\textwidth}{|l|X|}\hline
        $v, \, \some{v}, \, \none \in \Val$          & A value of any type; \ttf{some} and \none represent optional values.                             \\\hline
        $v_\bot \in \ValBot ::= v \, \vert \, \bot$  & A value that may be present $v$ or absent $\bot$.                                                \\\hline
        $v_d$                                        & A default value, predefined unique value for each type.                                          \\\hline
        $sty ::= ty \, \vert \, \mayB{ty}$           & A type, where \kwf{?} indicates an event type.                                                   \\\hline
        $\tyCtx[x] = sty$                            & A typing context \tyCtx mapping flow identifiers to their types.                                 \\\hline
        \mkcell{$\G[f] = \mathfoot{\component{f}{\mathtiny{x_i^+}}{\mathtiny{x_o^+}}{\mathtiny{eq^+}}{\mathtiny{x_m^+ = c^+}}}$                         \\
        $\G[F] = \service{F}{stmt^+}$}               & A program context \G mapping component and service names to their definitions.                   \\\hline
        $t \in \Time$                                & A timestamp.                                                                                     \\\hline
        $\ts \in \TStream{\Val}::= (t, v) \next \ts$ & A coinductive stream of increasing timestamped values.                                           \\\hline
        $\tck ::= (t, \true) \next \tck$             & A timestamped stream constantly equal to \true, named clock.                                     \\\hline
        $[\ts^+]$                                    & An implicitly finite list of streams $\ts^+$.                                                    \\\hline
        $\ts_1 \equiv \ts_2$                         & Equivalence of the timestamped streams $\ts_1$ and $\ts_2$.                                      \\\hline
        $K[x], I[x], O[x]$                           & Stream contexts $I$, $O$, and $K$ that map identifiers to their timestamped streams.             \\\hline
        \serviceDenot{\tck}{I}{F}{O}                 & The service $F$ is a function on \tck mapping context $I$ to context $O$.                        \\\hline
        \stmtDenot{\tck}{I}{K}{stmt}{O}              & The constraints of the statement $stmt$ are satisfied by the contexts $I$, $O$, and $K$ on \tck. \\\hline
        \sigmaDenot{\tck}{K}{\sigma}{\ts^+}          & The flow operation $\sigma$ is associated with the timestamped streams $\ts^+$ on \tck.          \\\hline
    \end{tabularx}
    \caption[Notations for the denotational semantics of \GRFRP.]%
    {Notations for the denotational semantics of \GRFRP.}
    \label{tab:semantics:grfrp:denot:notations}
\end{table}

We now introduce the notations used in this section to define the denotational semantics. These
notations, which build upon those introduced for the operational semantics in
\Cref{sec:semantics:grfrp:opsem}, are summarized in \Cref{tab:semantics:grfrp:denot:notations}.

In the denotational semantics of \GRFRP, flows are represented as infinite, ordered sequences of
timestamped values, called timestamped streams. These streams represent the values carried by the
flow at the times specified by the timestamps, similar to time-aware models such as the Tagged
Signal Model~\cite{DBLP:journals/tcad/LeeS98}. The set of all possible timestamps is denoted by
\Time, and a timestamp is denoted by $t$. Timestamped streams are elements of \TStream{\Val}, which
is the subset of \Stream{\Time \times \Val} containing only streams of increasing timestamps. They
are denoted by \ts and are coinductively constructed as $(t, v) \next \ts$. A signal (\resp event)
represented by $(t, v) \next \ts$ is updated (\resp occurs) at time $t$ with value $v$, and the
stream continues as \ts after $t$.
%
Clocks in this denotational semantics are timestamped streams of booleans, denoted by \tck.
They represent the execution times of computations and play a role similar to the logical clocks in
synchronous languages~\cite{DBLP:conf/emsoft/ColacoP03}, but are adapted to handle timestamps. A
stream \ts is said to be \emph{on} the clock \tck, written \on{\ts}{\tck}, when the
execution times of \tck make it possible to compute or use the values of \ts. In practice, this
means that the timestamps of \ts are included in the timestamps of \tck associated with the value
\true. This relationship is formally defined by the rules
\nameref{rule:semantics:streams:onclock:exact} and \nameref{rule:semantics:streams:onclock:pass}
(in \Cref{sec:semantics:streams:clock}). More information on timestamped clocks and timestamped
streams are provided in \Cref{sec:semantics:streams:def,sec:semantics:streams:clock}.
%
Finite lists of timestamped streams are denoted by $[\ts^+]$.
%
The contexts $I$ and $O$ map imported and exported flow identifiers ($x$) to their timestamped
streams ($I[x]$ or $O[x]$). Local flow identifiers ($x$) are stored in the context $K$ with their
timestamped streams ($K[x]$).
%
The concatenation of two timestamped stream contexts $X_1$ and $X_2$ ($X$ denoting either an input,
output, or local context $I$, $O$, or $K$), is defined as their disjoint union,
$X_1 \ctxconcat X_2$. A context mapping timestamped stream $\ts^+$ to identifiers $x^+$ is denoted
by $[(\ts/x)^+]$.
%
The equivalence of timestamped streams $\ts_1$ and $\ts_2$ is written $\ts_1 \equiv \ts_2$ and
states their timestamps and values are identical, as defined by the rule \nameref{eq:semantics:streams:equiv}
(in \Cref{sec:semantics:streams:def}).

Using the defined notations, we refine the purpose of the denotational semantics of \GRFRP. It
defines how to construct the timestamped streams corresponding to each identifier and flow
operation. It models flow statements as constraints on the contexts $I$, $O$, and $K$, and
represents services as functions from the input context $I$ to the output context $O$.%
\footnote{The determinism of the semantics, under causality of called components, is stated by
    \Cref{thm:semantics:grfrp:theorems:determinism} and proved in
    \Cref{sec:semantics:grfrp:theorems:determinism}.}
%
The following paragraphs explain in details the mutually recursive rules that define this semantics.

\paragraph{Services}

Consider a service \service{F}{stmt^+} defined in a program \G with typing context \tyCtx. Let \tck
be a timestamped clock. The denotational semantics of $F$ is a function that maps an input context
$I$ of timestamped streams to an output context $O$ of timestamped streams. The application of this
function is denoted by:
\begin{equation*}
    \serviceDenot{\tck}{I}{F}{O}.
\end{equation*}
The sets of valid input and output contexts are defined inductively over the structure of the
statements $stmt^+$ defining $F$.

First, we define the set of acceptable input contexts for a sequence of statements $stmt^+$, denoted
by \InputCtx{stmt^+}. The set is defined by induction on the statements, as shown in \Cref{%
    eq:semantics:grfrp:inputctx:import,%
    eq:semantics:grfrp:inputctx:export,%
    eq:semantics:grfrp:inputctx:decl%
}. An \kwf{import} statement adds a new identifier to the context, while other statements do not
affect the set of input contexts.

\begin{align}
    \InputCtx{\import{x}, stmt^{\prime+}}         & =
    \{ [\ts/x] \ctxconcat I \mid \ts \in \TStream{\Val}, I \in \InputCtx{stmt^{\prime+}} \} \label{eq:semantics:grfrp:inputctx:import} \\
    \InputCtx{\export{x}, stmt^{\prime+}}         & = \InputCtx{stmt^{\prime+}} \label{eq:semantics:grfrp:inputctx:export}             \\
    \InputCtx{\decl{x^+}{\sigma}, stmt^{\prime+}} & = \InputCtx{stmt^{\prime+}} \label{eq:semantics:grfrp:inputctx:decl}
\end{align}

Similarly, the set of acceptable output contexts, \OutputCtx{stmt^+}, contains mappings for all
identifiers introduced by \kwf{export} statements. The inductive definition in \Cref{%
    eq:semantics:grfrp:outputctx:import,%
    eq:semantics:grfrp:outputctx:export,%
    eq:semantics:grfrp:outputctx:decl%
} shows that only \kwf{export} statements extend the output context.

\begin{align}
    \OutputCtx{\import{x}, stmt^{\prime+}}         & = \OutputCtx{stmt^{\prime+}} \label{eq:semantics:grfrp:outputctx:import}            \\
    \OutputCtx{\export{x}, stmt^{\prime+}}         & =
    \{ [\ts/x] \ctxconcat O \mid \ts \in \TStream{\Val}, O \in \OutputCtx{stmt^{\prime+}} \} \label{eq:semantics:grfrp:outputctx:export} \\
    \OutputCtx{\decl{x^+}{\sigma}, stmt^{\prime+}} & = \OutputCtx{stmt^{\prime+}} \label{eq:semantics:grfrp:outputctx:decl}
\end{align}

Based on these context definitions, the denotational semantics of \service{F}{stmt^+}, denoted by
$\tyCtx, \G, \tck \vdashFrp F$, is of type $\InputCtx{stmt^+} \to \OutputCtx{stmt^+}$.
As established by the rule \nameref{rule:semantics:grfrp:denot:service}, for any input context
$I \in \InputCtx{stmt^+}$, the context $O \in \OutputCtx{stmt^+}$ is the output context if and only
if there exists a local context $K$ that satisfies the constraints imposed by the statements
$stmt^+$.

\begin{center}
    \begin{minipage}{0.9\textwidth}
        \centering
        \begin{mathpar}
            %%%%% Service
            \inferDENOT
            {
                \G[F] = \service{F}{stmt^+} \\
                \exists K, \, \left(\stmtDenot{\tck}{I}{K}{stmt}{O}\right)^+
            }
            {\serviceDenot{\tck}{I}{F}{O}}
            {Serv}
            \label{rule:semantics:grfrp:denot:service}
        \end{mathpar}
    \end{minipage}
\end{center}

The sets of valid local contexts, denoted by \LocalCtx{stmt^+}, is also defined inductively over the
structure of these statements. It contains mappings for all identifiers introduced by \kwf{import}
and local \kwf{let} statements. Indeed, the exported identifiers are already defined locally in the
service. As defined in \Cref{%
    eq:semantics:grfrp:localctx:import,%
    eq:semantics:grfrp:localctx:export,%
    eq:semantics:grfrp:localctx:decl%
}, only \kwf{import} and \kwf{let} statements extend the local context, but the \kwf{export}
statement constraints the local context.

\begin{align}
    \LocalCtx{\import{x}, stmt^{\prime+}}         & =
    \{ [\ts/x] \ctxconcat K \mid \ts \in \TStream{\Val}, K \in \LocalCtx{stmt^{\prime+}} \} \label{eq:semantics:grfrp:localctx:import} \\
    \LocalCtx{\export{x}, stmt^{\prime+}}         & =
    \{ K \mid K \in \LocalCtx{stmt^{\prime+}}, \: K[x] \in \TStream{\Val} \} \label{eq:semantics:grfrp:localctx:export}                \\
    \LocalCtx{\decl{x^+}{\sigma}, stmt^{\prime+}} & =
    \{ [(\ts/x)^+] \ctxconcat K \mid (\ts \in \TStream{\Val})^+, K \in \LocalCtx{stmt^{\prime+}} \} \label{eq:semantics:grfrp:localctx:decl}
\end{align}

We illustrate this semantics with an example.
For instance, consider the \aeb example from \Cref{fig:semantics:example:aeb}. In the execution
depicted in the chronogram from \Cref{fig:semantics:example:aeb:chronogram}, the denotational
semantics of the \idf{aeb} service is a function, defined on the timestamped clock \tck, that maps
the input context $I$ to the output context $O$:
\begin{equation*}
    \tck \vdashFrp \idf{aeb} [I] \Downarrow [O],
\end{equation*}
where the timestamped clock \tck, the timestamped streams corresponding to the identifiers
\idf{speed}, \idf{pedestrian\_l}, and \idf{pedestrian\_r} in the input context $I$, and the
timestamped stream corresponding to \idf{strat} in the output context $O$ are described below:
\begin{align*}
    \tck                              & = (t, \true) \next (t+1s, \true) \next (t+1.2s, \true) \next (t+1.5s, \true) \next (t+3s, \true) \cdots \\
    I[\idf{speed}]                    & = (t, 32.) \next (t+1.2s, 30.) \next (t+1.5s, 25.) \cdots                                               \\
    I[\mathfoot{\idf{pedestrian\_l}}] & = (t+1s, 20.) \cdots                                                                                    \\
    I[\mathfoot{\idf{pedestrian\_r}}] & = \cdots                                                                                                \\
    O[\idf{strat}]                    & = (t, \No) \next (t+1s, \Soft) \next (t+3s, \No) \cdots                                                 \\
\end{align*}

\paragraph{Flow statements}

A flow statement $stmt$ denotes constraints towards input, output, and local flow contexts $I$, $O$,
and $K$. The satisfaction of those constraints is written:
\begin{equation*}
    \stmtDenot{\tck}{I}{K}{stmt}{O},
\end{equation*}
and is defined by the following rules.

\subparagraph{Imports}

For an import statement \import{x}, rule \nameref{rule:semantics:grfrp:denot:stmt:import} specifies
that the local context $K$ associates the imported identifier $x$ with the same timestamped stream
as the input context $I$. In other words, it instantiates the input $x$ locally within the service.
Those identifiers are present in the contexts $I$ and $K$, as stated by
\Cref{eq:semantics:grfrp:inputctx:import,eq:semantics:grfrp:localctx:import}.
Furthermore, the condition \on{I[x]}{\tck} ensures that the timestamped stream $I[x]$ is on the
execution clock \tck, meaning its timestamps are included in those of \tck. This guarantees that
\tck is a correct execution clock for this import.

\begin{center}
    \begin{minipage}{0.9\textwidth}
        \centering
        \begin{mathpar}
            %%%%% Import
            \inferDENOTFlowStmt
            { I[x] \equiv K[x] \\ \on{I[x]}{\tck} }
            { \stmtDenot{\tck}{I}{K}{\import{x}}{O} }
            {Import}
            \label{rule:semantics:grfrp:denot:stmt:import}
        \end{mathpar}
    \end{minipage}
\end{center}

\subparagraph{Exports}

For an export statement \export{x}, rule \nameref{rule:semantics:grfrp:denot:stmt:export} specifies
that the local context $K$ associates the exported identifier $x$ with the same timestamped stream
as the output context $O$. In other words, it makes the local identifier $x$ available outside the
service.
Those identifiers are present in the contexts $O$ and $K$, as stated by
\Cref{eq:semantics:grfrp:outputctx:export,eq:semantics:grfrp:localctx:export}.
Furthermore, the condition \on{O[x]}{\tck} ensures that the timestamped stream $O[x]$ is on
the execution clock \tck, meaning its timestamps are included in those of \tck. This guarantees that
\tck is a correct execution clock for this export.

\begin{center}
    \begin{minipage}{0.9\textwidth}
        \centering
        \begin{mathpar}
            %%%%% Export
            \inferDENOTFlowStmt
            { O[x] \equiv K[x] \\ \on{O[x]}{\tck} }
            { \stmtDenot{\tck}{I}{K}{\export{x}}{O} }
            {Export}
            \label{rule:semantics:grfrp:denot:stmt:export}
        \end{mathpar}
    \end{minipage}
\end{center}

\subparagraph{Declarations}

And for a local declaration \decl{x^+}{\sigma} the rule
\nameref{rule:semantics:grfrp:denot:stmt:decl} states the local context $K$ associates to all
declared identifiers $x^+$ the timestamped stream corresponding to the evaluation of flow operation
$\sigma$ on the execution clock \tck.
Those identifiers are present in the local context $K$, as stated by
\Cref{eq:semantics:grfrp:localctx:decl}.

\begin{center}
    \begin{minipage}{0.9\textwidth}
        \centering
        \begin{mathpar}
            %%%%% Declaration
            \inferDENOTFlowStmt
            {
                \left(K[x] \equiv \ts\right)^+ \\
                \sigmaDenot{\tck}{K}{\sigma}{\ts^+}
            }
            { \stmtDenot{\tck}{I}{K}{\decl{x^+}{\sigma}}{O} }
            {Decl}
            \label{rule:semantics:grfrp:denot:stmt:decl}
        \end{mathpar}
    \end{minipage}
\end{center}

\paragraph{Flow operations}

A flow operation $\sigma$ is associated, in this denotational semantics, to a list of timestamped
streams $[\ts^+]$ that corresponds to its entire evaluation with a flow context $K$ and on an
execution clock \tck. This association is written:
\begin{equation*}
    \sigmaDenot{\tck}{K}{\sigma}{\ts^+},
\end{equation*}
and is defined by the following rules.

Most \GReact operations create timestamped streams using auxiliary functions that are defined
coinductively. We prove that these functions are well-defined using the coinduction
principle~\cite{DBLP:journals/mscs/KozenS17}, which ensures that the output timestamped streams are
constructed correctly and progress indefinitely. This principle, introduced earlier in
\Cref{sec:semantics:streams:coinduction}, provides a framework for reasoning about infinite
structures like timestamped streams. Specifically, it allows us to apply the coinduction hypothesis
only to streams that have \emph{progressed} and remain \emph{opaque}. For instance, to prove a
property on a timestamped stream $\ts = (t, v) \next \ts'$ ($\ts'$ has progressed), where no
assumptions are made about $\ts'$ ($\ts'$ is opaque), we invoke the coinduction hypothesis to
establish that the property holds for $\ts'$. Using this, we then prove that the property holds for
$(t, v) \next \ts'$ to conclude.

\subparagraph{Identifiers}

Identifiers ($x$) are associated with the corresponding timestamped stream in the local context
($K[x]$), see rule \nameref{rule:semantics:grfrp:denot:op:ident}. Moreover, this timestamped stream is on
the execution clock \tck, meaning that the clock \tck is a correct execution clock for this
identifier.

\begin{center}
    \begin{minipage}{0.9\textwidth}
        \centering
        \begin{mathpar}
            %%%%% Identifier
            \inferDENOTFlow
            {K[x] \equiv \ts \\ \on{K[x]}{\tck}}
            {\sigmaDenot{\tck}{K}{x}{\ts}}
            {ID}
            \label{rule:semantics:grfrp:denot:op:ident}
        \end{mathpar}
    \end{minipage}
\end{center}

\subparagraph{Component applications}

The denotational semantics of \GRSync, detailed in \Cref{sec:semantics:grsync:denot}, models flows
as simple streams (no timestamps) that have values on a simple clock $ck$ (no timestamps either).
This clock without timestamps represents a logical model of time, where each boolean value that is
\true (a tick) symbolize an instant of computation. In this semantics, see rule \nameref{rule:%
    semantics:grsync:denot:comp}, components are functions from input to output streams that are
on the same clock.

\begin{center}
    \begin{minipage}{0.95\textwidth}
        \centering
        \begin{mathpar}
            %%%%% Component application
            \inferDENOTFlow
            {
                \left(\sigmaDenot{\tck}{K}{\sigma}{\ts_i}\right)^+ \\
                \left(\on{\ts_i}{\tck}\right)^+ \\ \left(\on{\ts_o}{\tck}\right)^+ \\
                \G[f] = \component{f}{x_i^+}{x_o^+}{eq^+}{x_m^+ = c^+} \\
                \compDenot{\clock{\tck}}{f}{\clocked{\tyCtx[x_i]}{v_d}{\ts_i}{\tck}^+}{\clocked{\tyCtx[x_o]}{v_d}{\ts_o}{\tck}^+}
            }
            {\sigmaDenot{\tck}{K}{f(\sigma^+)}{\ts_o^+}}
            {App}
            \label{rule:semantics:grfrp:denot:op:app}
        \end{mathpar}
    \end{minipage}
\end{center}

The semantics of a component application $f(\sigma^+)$, see rule \nameref{rule:semantics:grfrp:%
    denot:op:app}, is a bridge between clocked streams of \GRSync and timestamped streams of \GRFRP.
%
From the timestamped clock \tck, we construct a simple clock \clock{\tck} that forgets the
timestamps as defined by \Cref{eq:semantics:grfrp:denot:aux:clock}, keeping only the logical symbol
of the instants of computation.
\begin{align}
    \clock{(t, \true) \next \tck} & = \true \step \clock{\tck} \label{eq:semantics:grfrp:denot:aux:clock}
\end{align}

The denotational semantics of \GRSync components is defined as a function that maps input simple
streams to output simple streams. These simple streams are obtained by projecting the input and
output timestamped streams onto a simple clock induced by \tck. This projection, denoted
\clocked{\tyCtx[x]}{v_d}{\ts}{\tck} for an identifier $x$ associated with a timestamped stream \ts,
is formally defined in
\Crefrange{eq:semantics:grfrp:denot:clocked:update}{eq:semantics:grfrp:denot:clocked:none}.
The projection uses \tyCtx to determine whether $x$ is a signal or an event, as the interpretation
of timestamped values depends on the flow type. For signals, the timestamped values represent
updates, with $v_d$ serving as the initial default value. The projection repeats the same value on
the simple clock ticks until the next update occurs
(\Crefrange{eq:semantics:grfrp:denot:clocked:update}{eq:semantics:grfrp:denot:clocked:repeat}).
For events, the timestamped values represent occurrences. The projection wraps these values in a
\ttf{some} option and inserts \none for all other ticks of the simple clock
(\Crefrange{eq:semantics:grfrp:denot:clocked:some}{eq:semantics:grfrp:denot:clocked:none}).

\begin{center}
    \scalebox{0.85}{
        \begin{minipage}{1.15\textwidth}
            \begin{align}
                \clocked{ty}{v^\explast}{(t, v) \next \ts}{(t', \true) \next \tck}        & = v \step \clocked{ty}{v}{\ts}{\tck}                                & \text{if} & \, t = t' \label{eq:semantics:grfrp:denot:clocked:update} \\
                \clocked{ty}{v^\explast}{(t, v) \next \ts}{(t', \true) \next \tck}        & = v^\explast \step \clocked{ty}{v^\explast}{(t, v) \next \ts}{\tck} & \text{if} & \, t > t' \label{eq:semantics:grfrp:denot:clocked:repeat} \\
                \clocked{\mayB{ty}}{v^\explast}{(t, v) \next \ts}{(t', \true) \next \tck} & = \some{v} \step \clocked{\mayB{ty}}{\some{v}}{\ts}{\tck}           & \text{if} & \, t = t' \label{eq:semantics:grfrp:denot:clocked:some}   \\
                \clocked{\mayB{ty}}{v^\explast}{(t, v) \next \ts}{(t', \true) \next \tck} & = \none \step \clocked{\mayB{ty}}{\none}{(t, v) \next \ts}{\tck}    & \text{if} & \, t > t' \label{eq:semantics:grfrp:denot:clocked:none}
            \end{align}
        \end{minipage}}
\end{center}

\begin{lemma}
    For all identifier $x$ typed by \tyCtx, all default value $v_d$, all timestamped clock \tck, and
    all timestamped stream \ts defined on \tck, the stream \clocked{\tyCtx[x]}{v_d}{\ts}{\tck} is
    well-defined.
\end{lemma}
\begin{skippableproof}
    Let $x$ be an identifier typed by \tyCtx, $v_d$ be a default value, \tck be a timestamped clock,
    and \ts be a timestamped stream defined on \tck.
    In all cases described by
    \Cref{%
        eq:semantics:grfrp:denot:clocked:update,%
        eq:semantics:grfrp:denot:clocked:repeat,%
        eq:semantics:grfrp:denot:clocked:some,%
        eq:semantics:grfrp:denot:clocked:none%
    }, the function advances at least one of the two input streams to its tail before its recursive
    call. As established in the analysis of coinductive stream pairs
    (\Cref{sec:semantics:streams:coinductive_constructs}), advancing even a single stream is
    sufficient to demonstrate progress for the pair as a whole, thereby justifying the use of the
    coinductive hypothesis.
    Moreover, the cases are exhaustive under the assumption that \ts is defined on \tck and that
    timestamps are strictly increasing, covering all possible combinations of timestamps, and
    ensuring opacity.
\end{skippableproof}

\subparagraph{Onchange}

In this denotational semantics, an operation \onchange{\sigma} is associated with the same timestamped
stream as its input signal $\sigma$, as described by the rule \nameref{rule:semantics:grfrp:denot:%
    op:onchange}. However, the meaning of this timestamped stream has changed: a timestamped value
for $\sigma$ denotes an update, while for the event \onchange{\sigma} it denotes an occurrence.

\begin{center}
    \begin{minipage}{0.9\textwidth}
        \centering
        \begin{mathpar}
            %%%%% OnChange
            \inferDENOTFlow
            {
                \sigmaDenot{\tck}{K}{\sigma}{\ts}
            }
            { \sigmaDenot{\tck}{K}{\onchange{\sigma}}{\ts} }
            {OChg}
            \label{rule:semantics:grfrp:denot:op:onchange}
        \end{mathpar}
    \end{minipage}
\end{center}

\subparagraph{Persist}

Most \GReact operations create timestamped streams using auxiliary functions.
%
The rule \nameref{rule:semantics:grfrp:denot:op:persist} states that, if $\sigma$ is an event
associated with the stream \ts, the semantics of \persist{\sigma} uses the function \persistAux{\ts}
{v^\explast}, defined in \Crefrange{eq:semantics:grfrp:denot:op:persist:aux:equal}{eq:semantics:grfrp:%
    denot:op:persist:aux:diff}, that creates a timestamped stream for the signal corresponding to
the persistence of $\sigma$.

\begin{center}
    \begin{minipage}{0.9\textwidth}
        \centering
        \begin{mathpar}
            %%%%% Persist
            \inferDENOTFlow
            {
                \sigmaDenot{\tck}{K}{\sigma}{\ts} \\
                \persistAux{\ts}{v_d} \equiv \ts'
            }
            { \sigmaDenot{\tck}{K}{\persist{\sigma}}{\ts'} }
            {Prst}
            \label{rule:semantics:grfrp:denot:op:persist}
        \end{mathpar}
    \end{minipage}
\end{center}

The function \persistAux{\ts}{v^\explast} filters the successive identical values by storing the last
occurrence value $v^\explast$. As a signal always has a value, all signals start with a default value
$v_d$, which is a predefined unique value for each type.
\begin{align}
    \persistAux{(t, v) \next \ts}{v^\explast} & =  \persistAux{\ts}{v^\explast}    & \text{if} & \, v = v^\explast \label{eq:semantics:grfrp:denot:op:persist:aux:equal}   \\
    \persistAux{(t, v) \next \ts}{v^\explast} & = (t, v) \next \persistAux{\ts}{v} & \text{if} & \, v \neq v^\explast \label{eq:semantics:grfrp:denot:op:persist:aux:diff}
\end{align}

\begin{lemma}
    For all timestamped stream \ts and all value $v^\explast \in \Val$, the timestamped stream
    \persistAux{\ts}{v^\explast} is well-defined.
\end{lemma}
\begin{skippableproof}
    Let \ts be a timestamped stream and $v^\explast \in \Val$.
    In all cases described by
    \Cref{%
        eq:semantics:grfrp:denot:op:persist:aux:equal,%
        eq:semantics:grfrp:denot:op:persist:aux:diff%
    }, the timestamped stream is consumed at each step, ensuring progress.
    Moreover, the cases are exhaustive ensuring opacity.
\end{skippableproof}

\subparagraph{Scan}

The rule \nameref{rule:semantics:grfrp:denot:op:scan} states that, if $\sigma$ is a signal
associated with the stream \ts, the semantics of the signal \scan{\sigma}{T} uses the function
\scanAux{\ts}{T}{t'}{v_{u\bot}}{v^\explast}, defined in \Crefrange{eq:semantics:grfrp:denot:op:scan:aux:1}
{eq:semantics:grfrp:denot:op:scan:aux:4}, that creates a timestamped stream corresponding to
$T$-periodical updates of $\sigma$.
%
As this operation has different timestamps compared to \ts, they have to be included in \tck.

\begin{center}
    \begin{minipage}{0.9\textwidth}
        \centering
        \begin{mathpar}
            %%%%% Scan
            \inferDENOTFlow
            {
                \sigmaDenot{\tck}{K}{\sigma}{\ts} \\
                \scanAux{\ts}{T}{T}{\bot}{v_d} \equiv \ts' \\
                \on{\ts'}{\tck}
            }
            { \sigmaDenot{\tck}{K}{\scan{\sigma}{T}}{\ts'} }
            {Scan}
            \label{rule:semantics:grfrp:denot:op:scan}
        \end{mathpar}
    \end{minipage}
\end{center}

The function \scanAux{\ts}{T}{t'}{v_{u\bot}}{v^\explast} stores the last update of \ts in $v_{u\bot}$ until
the next timer $t'$, when it produces this value:
if the next update of \ts occurs after $t'$ and there has been an update $v_u$, then it produces
$v_u$ at time $t'$ and continues for the next timer $t' + T$;
if the next update of \ts occurs after $t'$ but there has been no update ($v_{u\bot} = \bot$), then
no value is produced at time $t'$ and it continues for the next timer $t' + T$;
if the next update of \ts occurs before $t'$ and the value differs from the last periodic update
($v^\explast$), then the update is stored in $v_{u\bot}$; and
if the next update of \ts occurs before $t'$ but the value equals the last periodic update ($v^\explast$),
then $v_{u\bot}$ is reset to $\bot$.
\begin{align}
    \scanAux{(t, v) \next \ts}{T}{t'}{v_u}{v^\explast}       & = (t', v_u) \next \scanAux{(t, v) \next \ts}{T}{t'+T}{\bot}{v_u} & \text{if} & \, t > t' \label{eq:semantics:grfrp:denot:op:scan:aux:1}                            \\
    \scanAux{(t, v) \next \ts}{T}{t'}{\bot}{v^\explast}      & = \scanAux{(t, v) \next \ts}{T}{t'+T}{\bot}{v^\explast}          & \text{if} & \, t > t' \label{eq:semantics:grfrp:denot:op:scan:aux:2}                            \\
    \scanAux{(t, v) \next \ts}{T}{t'}{v_{u\bot}}{v^\explast} & = \scanAux{\ts}{T}{t'}{v}{v^\explast}                            & \text{if} & \, t \le t' \wedge v \neq v^\explast \label{eq:semantics:grfrp:denot:op:scan:aux:3} \\
    \scanAux{(t, v) \next \ts}{T}{t'}{v_{u\bot}}{v^\explast} & = \scanAux{\ts}{T}{t'}{\bot}{v^\explast}                         & \text{if} & \, t \le t' \wedge v = v^\explast \label{eq:semantics:grfrp:denot:op:scan:aux:4}
\end{align}

\begin{lemma}
    For all period $T$, all timestamped stream \ts, all timer $t'$, all bottom value
    $v_{u\bot} \in \ValBot$, and all value $v^\explast \in \Val$, the timestamped stream
    \scanAux{\ts}{T}{t'}{v_{u\bot}}{v^\explast} is well-defined.
\end{lemma}
\begin{skippableproof}
    The proof proceeds by coinduction. We consider the function \smf{scan^\#} as operating on a pair
    of streams: the explicit data stream \ts and an implicit timer stream, which we denote
    $\mathcal{T}$ and defined as follows, using the notations established in
    \Cref{sec:semantics:streams:coinductive_constructs}.
    \begin{equation*}
        \begin{split}
            \mathcal{T} & = t' + S_0^T                                \\
                        & = \texttt{trans}(t', S_0^T)                 \\
                        & = \texttt{trans}(t', 0 \step T \step S_2^T) \\
                        & = t' \step \texttt{trans}(t' + T, S_1^T)    \\
                        & = t' \step (t'+T + S_1^T)
        \end{split}
    \end{equation*}
    %
    To prove the function is well-defined, we must show that for any state, it makes progress on
    this pair before any recursive call. We analyze the four defining cases from
    \Cref{%
        eq:semantics:grfrp:denot:op:scan:aux:1,%
        eq:semantics:grfrp:denot:op:scan:aux:2,%
        eq:semantics:grfrp:denot:op:scan:aux:3,%
        eq:semantics:grfrp:denot:op:scan:aux:4%
    }:
    %
    \begin{description}
        \item[Cases 1 and 2 ($t > t'$)] The recursive call is made with the timer parameter advanced
              to $t' + T$. This corresponds to consuming the head of the timer stream $\mathcal{T}$
              (which is $t'$) and continuing with its tail, $t'+T + S_1^T$. The data stream \ts is
              not advanced. Thus, the pair $(\ts, \mathcal{T})$ makes progress.

        \item[Cases 3 and 4 ($t \le t'$)] The recursive call is made with the data stream advanced
              to its tail, \ts. The timer parameter $t'$ is unchanged, meaning the timer stream
              $\mathcal{T}$ is not advanced. Thus, the pair $(\ts, \mathcal{T})$ makes progress.
    \end{description}
    %
    In all four exhaustive cases, the function advances at least one of the two component streams of
    the pair before the recursive call. As established in the analysis of coinductive
    stream pairs (\Cref{sec:semantics:streams:coinductive_constructs}), advancing even a single
    stream is sufficient to demonstrate progress for the pair as a whole.
    %
    Because progress is established, we can apply the \emph{coinductive hypothesis} and assume that
    any recursive call to \smf{scan^\#} yields a well-defined timestamped stream. Since the
    function's output is either constructed with a new head (Case 1) or is the direct result of
    the recursive call, we conclude that \scanAux{\ts}{T}{t'}{v_{u\bot}}{v^\explast} is
    well-defined.
\end{skippableproof}

\subparagraph{Sample}

The rule \nameref{rule:semantics:grfrp:denot:op:sample} states that, if $\sigma$ is an event
associated with the stream \ts, the semantics of the event \sample{\sigma}{T} uses the function
\sampleAux{\ts}{T}{t'}{v_{u\bot}}, defined in \Crefrange{eq:semantics:grfrp:denot:op:sample:aux:1}
{eq:semantics:grfrp:denot:op:sample:aux:3}, that creates a timestamped stream corresponding to
$T$-periodical occurrences of $\sigma$.
%
As this operation has different timestamps compared to \ts, they have to be included in \tck.

\begin{center}
    \begin{minipage}{0.9\textwidth}
        \centering
        \begin{mathpar}
            %%%%% Sample
            \inferDENOTFlow
            {
                \sigmaDenot{\tck}{K}{\sigma}{\ts} \\
                \sampleAux{\ts}{T}{T}{\bot} \equiv \ts' \\
                \on{\ts'}{\tck}
            }
            { \sigmaDenot{\tck}{K}{\sample{\sigma}{T}}{\ts'}}
            {Samp}
            \label{rule:semantics:grfrp:denot:op:sample}
        \end{mathpar}
    \end{minipage}
\end{center}

The function \sampleAux{\ts}{T}{t'}{v_{u\bot}} stores the last occurrence of \ts in $v_{o\bot}$ until
the next period $t'$, when it produces this value:
if the next occurrence of \ts occurs after $t'$ and there has been an occurrence $v_u$, then it
produces $v_o$ at time $t'$ and continues for the next period $t' + T$;
if the next occurrence of \ts occurs after $t'$ but there has been not occurrence ($v_{o\bot} =
    \bot$), then no value is produced at time $t'$ and it continues for the next period $t' + T$;
and
if the next occurrence of \ts occurs before $t'$, then the occurrence is stored in $v_{o\bot}$.
\begin{align}
    \sampleAux{(t, v) \next \ts}{T}{t'}{v_o}       & = (t', v_o) \next \sampleAux{(t, v) \next \ts}{T}{t'+T}{\bot} & \text{if} & \, t > t' \label{eq:semantics:grfrp:denot:op:sample:aux:1}   \\
    \sampleAux{(t, v) \next \ts}{T}{t'}{\bot}      & = \sampleAux{(t, v) \next \ts}{T}{t'+T}{\bot}                 & \text{if} & \,  t > t' \label{eq:semantics:grfrp:denot:op:sample:aux:2}  \\
    \sampleAux{(t, v) \next \ts}{T}{t'}{v_{o\bot}} & = \sampleAux{\ts}{T}{t'}{v}                                   & \text{if} & \, t \le t' \label{eq:semantics:grfrp:denot:op:sample:aux:3}
\end{align}

\begin{lemma}
    For all period $T$, all timestamped stream \ts, all time $t'$, and all bottom value
    $v_{u\bot} \in \ValBot$, the timestamped stream \sampleAux{\ts}{T}{t'}{v_{u\bot}} is well-defined.
\end{lemma}
\begin{skippableproof}
    The proof follows the same coinductive reasoning established for the \smf{scan^\#} function. We
    consider \smf{sample^\#} as a function operating on a pair of streams: the explicit data stream
    \ts and an implicit timer stream, $\mathcal{T} = t' + S_0^T$.
    %
    To prove the function is well-defined, we show that it makes progress on this pair in all
    defining cases from \Cref{%
        eq:semantics:grfrp:denot:op:sample:aux:1,%
        eq:semantics:grfrp:denot:op:sample:aux:2,%
        eq:semantics:grfrp:denot:op:sample:aux:3%
    }:
    %
    \begin{description}
        \item[Cases 1 and 2 ($t > t'$)] The timer parameter is advanced to $t' + T$. This
              corresponds to consuming the head of the timer stream $\mathcal{T}$ while the data
              stream \ts is not advanced. The pair thus makes progress.

        \item[Case 3 ($t \le t'$)] The data stream is advanced to its tail, \ts, while the timer
              parameter $t'$ is unchanged. The pair thus makes progress.
    \end{description}
    %
    In all three exhaustive cases, the function advances at least one of the two component streams
    of the pair before the recursive call. As established in
    \Cref{sec:semantics:streams:coinductive_constructs}, this is sufficient to demonstrate progress
    for the pair as a whole.
    %
    Because progress is established, the coinductive hypothesis allows us to assume that any
    recursive call to \smf{sample^\#} yields a well-defined timestamped stream. Since the output of
    the function is either constructed with a new head (Case 1) or is the direct result of the
    recursive call, we conclude that \sampleAux{\ts}{T}{t'}{v_{o\bot}} is well-defined.
\end{skippableproof}

\subparagraph{Timeout}

The rule \nameref{rule:semantics:grfrp:denot:op:timeout} states that, if $\sigma$ is an event
associated with the stream \ts, the semantics of the event \timeout{\sigma}{T} uses the function
\timeoutAux{\ts}{T}{t'}, defined in \Crefrange{eq:semantics:grfrp:denot:op:timeout:aux:1}
{eq:semantics:grfrp:denot:op:timeout:aux:2}, that creates a timestamped stream with unit values when
$\sigma$ has not occurred for $T$ milliseconds.
%
As this operation has different timestamps compared to \ts, they have to be included in \tck.

\begin{center}
    \begin{minipage}{0.9\textwidth}
        \centering
        \begin{mathpar}
            %%%%% Timeout
            \inferDENOTFlow
            {
                \sigmaDenot{\tck}{K}{\sigma}{\ts} \\
                \timeoutAux{\ts}{T}{T} \equiv \ts' \\
                \on{\ts'}{\tck}
            }
            { \sigmaDenot{\tck}{K}{\timeout{\sigma}{T}}{\ts'}}
            {TOut}
            \label{rule:semantics:grfrp:denot:op:timeout}
        \end{mathpar}
    \end{minipage}
\end{center}

If the next occurrence of \ts occurs after the timer $t'$, then \timeoutAux{\ts}{T}{t'} produces a
unit value $()$ at time $t'$ and continues with the next timer set to $t' + T$. Otherwise, it sets
the next timer to $t + T$, where $t$ is the timestamp of the occurrence of \ts.
\begin{align}
    \timeoutAux{(t, v) \next \ts}{T}{t'} & = (t', ()) \next \timeoutAux{(t, v) \next \ts}{T}{t' + T} & \text{if} & \, t > t' \label{eq:semantics:grfrp:denot:op:timeout:aux:1}   \\
    \timeoutAux{(t, v) \next \ts}{T}{t'} & = \timeoutAux{\ts}{T}{t + T}                              & \text{if} & \, t \le t' \label{eq:semantics:grfrp:denot:op:timeout:aux:2}
\end{align}

\begin{lemma}
    For all period $T$, all timestamped stream \ts, and all timer $t'$, the timestamped stream
    \timeoutAux{\ts}{T}{t'} is well-defined.
\end{lemma}
\begin{skippableproof}
    The proof proceeds by coinduction, following the same structure used for \smf{scan^\#} and
    \smf{sample^\#}. We consider \smf{timeout^\#} as a function operating on a pair of streams: the
    explicit data stream \ts and an implicit timer stream, which we denote $\mathcal{T}$. Unlike the
    previous functions, the timer stream here is not independent of the data stream. At any point,
    $\mathcal{T}$ is defined by its current starting time, say $t'$, as $\mathcal{T} = t' + S_0^T$.
    %
    To prove the function is well-defined, we must show that it makes progress on this pair in
    all defining cases from \Cref{%
        eq:semantics:grfrp:denot:op:timeout:aux:1,%
        eq:semantics:grfrp:denot:op:timeout:aux:2%
    }:
    %
    \begin{description}
        \item[Case 1 ($t > t'$)] The recursive call is made with the timer parameter advanced to
              $t' + T$. As shown in the proof for \smf{scan^\#}, this corresponds to consuming
              the head of the current timer stream $\mathcal{T} = t' + S_0^T$ and continuing with
              its tail $t'+T + S_1^T$. The data stream \ts is not advanced. The pair thus makes
              progress.

        \item[Case 2 ($t \le t'$)] The data stream is advanced to its tail, \ts. Simultaneously,
              the timer parameter is reset to $t + T$. This redefines the implicit timer stream for
              the subsequent computation to $\mathcal{T}' = t + S_0^T = t \step ((t+T) + S_1^T)$,
              consumes its head, and make the recursive call on its tail. Since both the data
              stream and the timer stream have progressed (one to its tail, the other to the tail of
              its new definition), the pair as a whole makes progress.
    \end{description}
    %
    In both exhaustive cases, the function makes progress on the stream pair before the recursive
    call. Because progress is established, we can apply the \emph{coinductive hypothesis} and assume
    that any recursive call to \smf{timeout^\#} yields a well-defined timestamped stream. Since the
    output of the function is either constructed with a new head (Case 1) or is the direct result of
    the recursive call, we conclude that \smf{timeout^\#} is well-defined for any valid initial
    timer $t'$.
\end{skippableproof}

\subparagraph{Merge}

The rule \nameref{rule:semantics:grfrp:denot:op:merge} states that, if $\sigma_1$ and $\sigma_2$ are
two events associated with the streams $\ts_1$ an $\ts_2$ respectively, the semantics of the event
\mergeop{\sigma_1}{\sigma_2} uses the function \mergeAux{\ts_1}{\ts_2}, defined in \Crefrange{eq:%
    semantics:grfrp:denot:op:merge:aux:1}{eq:semantics:grfrp:denot:op:merge:aux:3}, that creates a
timestamped stream that contains the union of both occurrences, with priority to the first event
$\sigma_1$ in case they occur simultaneously.

\begin{center}
    \begin{minipage}{0.9\textwidth}
        \centering
        \begin{mathpar}
            %%%%% Merge
            \inferDENOTFlow
            {
                \sigmaDenot{\tck}{K}{\sigma_1}{\ts_1} \\
                \sigmaDenot{\tck}{K}{\sigma_2}{\ts_2} \\
                \mergeAux{\ts_1}{\ts_2} \equiv \ts'
            }
            { \sigmaDenot{\tck}{K}{\mergeop{\sigma_1}{\sigma_2}}{\ts'}}
            {Mrg}
            \label{rule:semantics:grfrp:denot:op:merge}
        \end{mathpar}
    \end{minipage}
\end{center}

If the next occurrence of $\ts_1$ occurs at the same time as $\ts_2$, then the value of $\ts_1$ is emitted
by \mergeAux{\ts_1}{\ts_2} and the one of $\ts_2$ is dropped; if one occurs before the other, its value
is emitted.
\begin{align}
    \mergeAux{(t_1, v_1) \next \ts_1}{(t_2, v_2) \next \ts_2} & = (t_1, v_1) \next \mergeAux{\ts_1}{\ts_2}                   &
    \text{if}                                                 & \, t_1 = t_2 \label{eq:semantics:grfrp:denot:op:merge:aux:1}   \\
    \mergeAux{(t_1, v_1) \next \ts_1}{(t_2, v_2) \next \ts_2} & = (t_1, v_1) \next \mergeAux{\ts_1}{(t_2, v_2) \next \ts_2}  &
    \text{if}                                                 & \, t_1 < t_2 \label{eq:semantics:grfrp:denot:op:merge:aux:2}   \\
    \mergeAux{(t_1, v_1) \next \ts_1}{(t_2, v_2) \next \ts_2} & = (t_2, v_2) \next \mergeAux{(t_1, v_1) \next \ts_1}{\ts_2}  &
    \text{if}                                                 & \, t_1 > t_2 \label{eq:semantics:grfrp:denot:op:merge:aux:3}
\end{align}

\begin{lemma}
    For all timestamped streams $\ts_1$ and $\ts_2$, the timestamped stream \mergeAux{\ts_1}{\ts_2}
    is well-defined.
\end{lemma}
\begin{skippableproof}
    The proof proceeds by coinduction. We consider \smf{merge^\#} as a function operating on the
    coinductive stream pair $(\ts_1, \ts_2)$. To prove the function is well-defined, we must show
    that it makes progress on this pair before any recursive call. We analyze the three defining
    cases from \Cref{%
        eq:semantics:grfrp:denot:op:merge:aux:1,%
        eq:semantics:grfrp:denot:op:merge:aux:2,%
        eq:semantics:grfrp:denot:op:merge:aux:3%
    }:
    %
    \begin{description}
        \item[Case 1 ($t_1 < t_2$)] The recursive call is made on $(\texttt{tail}(\ts_1), \ts_2)$.
              The first component of the pair is advanced. The pair thus makes progress.

        \item[Case 2 ($t_1 > t_2$)] The recursive call is made on $(\ts_1, \texttt{tail}(\ts_2))$.
              The second component of the pair is advanced. The pair thus makes progress.

        \item[Case 3 ($t_1 = t_2$)] The recursive call is made on $(\texttt{tail}(\ts_1), \texttt{tail}(\ts_2))$.
              Both components of the pair are advanced. The pair thus makes progress.
    \end{description}
    %
    In all three exhaustive cases, the function advances at least one of the two component streams
    before the recursive call. As established in the analysis of coinductive stream pairs
    (\Cref{sec:semantics:streams:coinductive_constructs}), advancing even a single component stream
    is sufficient to demonstrate progress for the pair as a whole.
    %
    Because progress is established, we can apply the \emph{coinductive hypothesis} and assume that
    any recursive call to \smf{merge^\#} yields a well-defined timestamped stream. Since the
    output of the function is always constructed by prepending a new head to the result of the
    recursive call, we conclude that \mergeAux{\ts_1}{\ts_2} is well-defined.
\end{skippableproof}

\subparagraph{Throttle}

The rule \nameref{rule:semantics:grfrp:denot:op:throttle} states that, if $\sigma$ is a signal
associated with the stream \ts, the semantics of the signal \throttle{\sigma}{\Delta} uses the
function \throttleAux{\ts}{\Delta}{v^\explast}, defined in \Crefrange{eq:semantics:grfrp:denot:op:throttle:%
    aux:1}{eq:semantics:grfrp:denot:op:throttle:aux:2}, that creates a timestamped stream of the
updates from $\sigma$ that have a difference greater that $\Delta$ from each other.

\begin{center}
    \begin{minipage}{0.9\textwidth}
        \centering
        \begin{mathpar}
            %%%%% Throttle
            \inferDENOTFlow
            {
                \sigmaDenot{\tck}{K}{\sigma}{\ts} \\
                \throttleAux{\ts}{\Delta}{v_d} \equiv \ts'
            }
            { \sigmaDenot{\tck}{K}{\throttle{\sigma}{\Delta}}{\ts'}}
            {Thrt}
            \label{rule:semantics:grfrp:denot:op:throttle}
        \end{mathpar}
    \end{minipage}
\end{center}

The function \throttleAux{\ts}{\Delta}{v^\explast} stores its last update in $v^\explast$ to filter the changes that
are not significant:
if the difference between the next update of \ts and the last value $v^\explast$ is greater that $\Delta$,
then this value passes and $v^\explast$ is updated; otherwise, the value is filtered.
\begin{align}
    \throttleAux{(t', v') \next \ts'}{\Delta}{v^\explast} & = (t', v') \next \throttleAux{\ts'}{\Delta}{v'}                                    &
    \text{if}                                             & \, |v' - v^\explast| \ge \Delta \label{eq:semantics:grfrp:denot:op:throttle:aux:1}   \\
    \throttleAux{(t', v') \next \ts'}{\Delta}{v^\explast} & = \throttleAux{\ts'}{\Delta}{v^\explast}                                           &
    \text{if}                                             & \, |v' - v^\explast| > \Delta \label{eq:semantics:grfrp:denot:op:throttle:aux:2}
\end{align}

\begin{lemma}
    For all value $\Delta \in \Val$, all timestamped stream \ts, and all value $v^\explast \in \Val$,
    the timestamped stream \throttleAux{\ts}{\Delta}{v^\explast} is well-defined.
\end{lemma}
\begin{skippableproof}
    Let $\Delta \in \Val$, \ts be a timestamped stream, and $v^\explast \in \Val$.
    In all cases described by
    \Cref{%
        eq:semantics:grfrp:denot:op:throttle:aux:1,%
        eq:semantics:grfrp:denot:op:throttle:aux:2%
    }, \ts is consumed at each step, ensuring progress. Moreover, the cases are exhaustive ensuring
    opacity.
\end{skippableproof}

\setcounter{subparagraph}{0}
\subparagraph{Time}

The operation \timeop returns the timestamps of the clock \tck. This behavior is defined
by rule \nameref{rule:semantics:grfrp:denot:op:time}, which relies on the auxiliary function
\timeAux{\tck}. This function, defined in \Cref{eq:semantics:grfrp:denot:op:time:aux}, creates the
corresponding timestamped stream from \tck.

\begin{center}
    \begin{minipage}{0.9\textwidth}
        \centering
        \begin{mathpar}
            %%%%% Time
            \inferDENOTFlow
            { \timeAux{\tck} \equiv \ts }
            { \sigmaDenot{\tck}{K}{\timeop}{\ts}}
            {Time}
            \label{rule:semantics:grfrp:denot:op:time}
        \end{mathpar}
    \end{minipage}
\end{center}

\begin{align}
    \timeAux{(t, \true) \next \tck} & = (t, t) \next \timeAux{\tck} \label{eq:semantics:grfrp:denot:op:time:aux}
\end{align}

\begin{lemma}
    For all timestamped clock \tck, the timestamped stream \timeAux{\tck} is well-defined.
\end{lemma}
\begin{skippableproof}
    In \Cref{eq:semantics:grfrp:denot:op:time:aux}, the clock is consumed without making any
    assumptions about the remaining tail of the clock, thereby ensuring both progress and opacity.
\end{skippableproof}

\paragraph{Summary}
This section presented the denotational semantics for \GRFRP, which interprets flows as infinite
timestamped streams. This model provides a declarative, time-aware perspective on system behavior.
Services are modeled as functions that map input stream contexts to output stream contexts, while
statements are interpreted as constraints on these contexts. The semantics of each flow operation is
defined by a coinductive stream-transforming function. A key aspect of this formalization is the
semantic bridge to \GRSync: component applications are handled by projecting timestamped streams
onto a logical clock, allowing the synchronous, clocked semantics of components to be integrated
into the asynchronous, timestamped world of services. This provides a functional interpretation of
\GRFRP programs, abstracting away from execution order to focus on the temporal relationships
between dataflows.
